% ===================================================================
%  نقشہ نمبر 5 — گھر تے اوہدا بݨنا۔
%  Traduction 2026 d'apres texto/io, controlee sur texto/fr, texto/en
%  et texto/hi. Consigne : voir texto/pnb/10-naqsha-01.tex.
%
%  Le tableau 5 est un DIALOGUE, et les deux volumes composent le nom
%  du parleur en petites capitales : on garde \textsc{}, que le rendu
%  laisse tomber en gardant le texte.
%
%  Les renvois du plan sont des LETTRES — (a) le couloir, (b) le
%  salon... — et non des numeros. On les ecrit comme l'ido les ecrit,
%  y compris la ou l'imprimeur a compose « (1) » pour « (l) ».
%
%  LES DEUX NOMS SONT TRANSCRITS, non traduits. Le shahmoukhi rend
%  les noms francais a la lettre, et il l'a fait pour tous les
%  prenoms de ces seize tableaux : M. Leblanc est « مسٹر لبلاں »,
%  M. Roux est « مسٹر رُو ». Le hindi, lui, traduisait la couleur du
%  premier et transcrivait le second ; on ne l'a pas suivi, parce
%  qu'une colonne doit se tenir a un seul parti.
% ===================================================================

%%K t05-apar-1 apar
\VUsaut{6.0mm}
\VUtitre{15.0pt}{60}{نقشہ نمبر 5}
\VUsaut{1.4mm}
\VUtitre{11.4pt}{0}{گھر تے اوہدا بݨنا۔}
\VUsaut{2.2mm}

%%K t05-tit-1 sub
\VUpk{10.2pt}{40}{چنگی ملاقات۔}
\VUsaut{3.0mm}

%%K t05-01-1 p
\textsc{جان}۔ --- چاچا جی، میں بہت جاݨنا چاہنا واں پئی \VUgras{گھر} کِویں بݨدا
اے۔

%%K t05-01-2 p
\textsc{چاچا} \textsuperscript{(80)}۔ --- ایہہ سوکھا اے، پُتر؛ میرے نال آ تے میں
تینوں اوہ گھر وکھاواں گا جیہڑا مسٹر برنارڈ \textsuperscript{(79)} بݨوا رہے نیں تے
جیدے وچ میں رہݨ جا رہیا واں۔

%%K t05-01-3 p
\textsc{جان}۔ --- تے ایس گھر دا \VUgras{نقشہ} \textsuperscript{(76)} کیہنے بݨایا؟

%%K t05-01-4 p
\textsc{چاچا}۔ --- \VUgras{معمار} \textsuperscript{(75)} نے؛ اوہناں بہت صاف نقشہ
بݨایا، اوہ من لیا گیا، تے \VUgras{ٹھیکیدار} \textsuperscript{(77)} نے کم شروع کر
دِتا۔

%%K t05-01-5 p
\textsc{جان}۔ --- تے ٹھیکیدار کوݨ نیں؟

%%K t05-01-6 p
\textsc{چاچا}۔ --- مسٹر لبلاں، تیرے ساتھی جیمز دے پیو۔ اوہ معمار مسٹر رُو نال رل
کے مزدوراں نوں ہدایتاں دیندے نیں۔

%%K t05-01-7 p
لے! مسٹر لبلاں ٹھیک ویلے اپݨا \VUgras{پیمانہ} \textsuperscript{(78)} چُکی آ رہے
نیں، تے مسٹر رُو اپݨا نقشہ چُکی۔

%%K t05-01-8 p
سلام، جناب۔ تُسیں کِویں او؟

%%K t05-01-9 p
\textsc{مسٹر رُو}۔ --- بہت چنگے، شکریہ۔ سلام جان؛ مُنڈا کِنّا وڈا ہو گیا اے!

%%K t05-01-10 p
\textsc{چاچا}۔ --- سچی گل، ایہہ پورا نِکا بندہ اے۔ ایہہ بہت سنجیدہ اے؛ جیہڑا کجھ
ویکھدا اے، اوہدی وضاحت کرنی پیندی اے۔ اج ایہہ جاݨنا چاہندا اے پئی گھر کِویں
بݨدا اے۔

%%K t05-tit-2 sub
\VUpk{10.2pt}{40}{مزدور۔}
\VUsaut{3.0mm}

%%K t05-01-11 p
\textsc{مسٹر لبلاں}۔ --- تے میرے نال آ، میرے نِکے یار، تے میں تینوں ہُݨے سمجھا
دینا واں، کیوں جے مینوں سِکھݨ دے چاہواݨ نِکے بہت چنگے لگدے نیں۔

%%K t05-01-12 p
اوہ مزدور ویکھنا ایں، جیدے ہتھ وچ \VUgras{کہی} \textsuperscript{(82)} اے: اوہ
\VUgras{کھدائی کرن والا} \textsuperscript{(81)} اے۔ اوہ پاݨی دی نالی وِچھاؤݨ لئی اک
\VUgras{کھال} \textsuperscript{(46)} پُٹ رہیا اے۔ اوسے نے اوہ زمین وی پُٹی سی
جِتھے گھر کھڑا اے، تاں جو راج چوہاں \VUgras{کندھاں} چُک سکے۔ اوہناں کندھاں دا
جیہڑا حصہ زمین وچ اے اوہنوں \VUgras{نیہہ} \textsuperscript{(2)} آکھدے نیں۔ نیہہ
نال گھِری تھاں \VUgras{تہہ خانہ} \textsuperscript{(3)} بݨاندی اے۔ ایس تہہ خانے وچ
اک نِکی کھڑکی اے جیہنوں \VUgras{روشندان} \textsuperscript{(5)} آکھدے نیں، اک
\VUgras{دروازہ} \textsuperscript{(4)} تے اک پوڑی۔

%%K t05-01-13 p
\VUgras{راج} \textsuperscript{(108)} کندھاں چُکدا گیا، \VUgras{بوہیاں}
\textsuperscript{(8)} تے \VUgras{کھڑکیاں} \textsuperscript{(17)} لئی تھاں چھڈدا
ہویا۔

%%K t05-01-14 p
\textsc{جان}۔ --- پر مینوں تے ایہہ راج وکھائی ای نئیں دیندا!

%%K t05-01-15 p
\textsc{مسٹر لبلاں}۔ --- اوہدا گھر دا کم ہُݨ پورا ہو چُکیا اے۔ اوہ اودھر
\VUgras{طویلے} \textsuperscript{(54)} کول اے، اپݨے \VUgras{پاڑ} \textsuperscript{(110)}
اُتے چڑھیا، \VUgras{دھوبی گھر} \textsuperscript{(56)} دی کندھ بݨا رہیا اے۔

%%K t05-01-16 p
اوہدے کول اک کرنی، اک نانْد تے اک \VUgras{ساہل} \textsuperscript{(109)} اے۔ پر ایہہ
ناں تینوں یاد نہ رہݨ گے۔

%%K t05-01-17 p
\textsc{جان}۔ --- کیوں نہ رہݨ گے، میں چاچا جی کولوں اک نِکی کرنی تے اک نانْد
منگاں گا، تے ساہل تے میں سُتلی دے ٹوٹے نال آپ ای بݨا لواں گا۔

%%K t05-tit-3 sub
\VUsaut{3.0mm}
\VUpk{10.2pt}{40}{سامان۔}
\VUsaut{3.0mm}

%%K t05-01-18 p
\textsc{چاچا}۔ --- واہ، بہت چنگا۔ پر دس تے جان، گھر بݨاؤݨ لئی کیہ چاہیدا اے؟

%%K t05-01-19 p
\textsc{جان}۔ --- \VUgras{تراشیا ہویا پتھر} \textsuperscript{(59)} تے
\VUgras{گارا} \textsuperscript{(65)} چاہیدا اے۔

%%K t05-01-20 p
\textsc{چاچا}۔ --- ٹھیک۔ اوس وڈی ہیٹ والے بندے نوں ویکھ، اپݨے \VUgras{گڈے}
\textsuperscript{(136)} اُتے۔ اوہ \VUgras{گڈے والا} \textsuperscript{(135)} اے۔ اوہ
ہر دن ایتھے \VUgras{پتھر دے ڈھیمے} \textsuperscript{(60)} لیاندا اے۔ اوہ اپݨا گڈا
وِہڑے وچ خالی کردا اے، تے \VUgras{سنگ تراش} \textsuperscript{(83)} ڈھیمیاں نوں
تراشدے تے اپݨی \VUgras{پتھر آری} \textsuperscript{(85)} نال چیردے نیں۔ اک
\VUgras{مزدور} \textsuperscript{(146)} اوہناں نوں راج تیکر اپڑاندا اے، جیہڑا اوہناں
نوں اپݨی تھاں بہاندا اے۔

%%K t05-01-21 p
ایس دوران \VUgras{گارا بݨاؤݨ والا} \textsuperscript{(87)} گارا تیار کردا اے۔
اوہنوں \VUgras{ریت} \textsuperscript{(62)}، \VUgras{چوݨا} \textsuperscript{(63)} تے
\VUgras{پاݨی} \textsuperscript{(64)} چاہیدا اے۔ چوݨا اوہدے کول اک \VUgras{بوری}
\textsuperscript{(71)} وچ اے؛ اک \VUgras{راج دا مددگار} \textsuperscript{(111)}
اوہنوں \VUgras{ٹھیلے} \textsuperscript{(112)} وچ ریت لیاندا اے، تے
\VUgras{شاگرد} \textsuperscript{(113)} پمپ \textsuperscript{(58)} اُتے اے۔ اوہ اک
\VUgras{بالٹی} \textsuperscript{(114)} بھر رہیا اے۔ گارا چھیتی ای تیار ہو جاوے گا۔
\VUgras{گارا چُکݨ والا} \textsuperscript{(107)} اوہنوں لے کے پوڑی توں پاڑ اُتے چڑھا
لے جاندا اے، جِتھے اک راج گھر دا اوہ پاسا پورا کر رہیا اے۔

%%K t05-01-22 p
تے ہُݨ دس، جیہڑا مزدور \VUgras{چھت} \textsuperscript{(31)} بݨاندا اے اوہنوں کیہ
آکھدے نیں؟

%%K t05-01-23 p
\textsc{جان}۔ --- اوہ \VUgras{چھت بݨاؤݨ والا} \textsuperscript{(118)} اے۔

%%K t05-tit-4 sub
\VUsaut{3.0mm}
\VUpk{10.2pt}{40}{گھر دی چھت۔}
\VUsaut{3.0mm}

%%K t05-01-24 p
\textsc{مسٹر لبلاں}۔ --- ہاں، پر پہلاں \VUgras{ترکھاݨ} \textsuperscript{(115)}
آندا اے۔

%%K t05-01-25 p
ترکھاݨ \VUgras{چھت دا ڈھانچہ} \textsuperscript{(35)} بݨاندا اے۔ اوہ
\VUgras{کڑیاں} \textsuperscript{(37)} نوں \VUgras{کِلیاں} \textsuperscript{(117)}
نال جوڑدا اے۔ اُتے اوہ مزدور، اپݨی چوڑی مخمل دی نِکر وچ، ترکھاݨ نیں۔ اک نِکی کڑی
پھڑی ہوئی اے تے دوجا اپݨے \VUgras{کہاڑے} \textsuperscript{(116)} نال اک کِلا گھڑ
رہیا اے۔ \VUgras{چمنی} \textsuperscript{(41)} کول والا چھت بݨاؤݨ والا اے۔ اوہ
\VUgras{شہتیراں} (*) اُتے پٹیاں ٹھوکدا اے، تے پٹیاں اُتے \VUgras{سلیٹاں}
\textsuperscript{(66)}۔ کدے کدے اوہ کھپریلاں وِچھاندا اے۔

%%K t05-noto-1 noto
\VUnotes{143.2mm}{%
(*) \VUgras{Balk-o} = جرمن \textit{Balken}، انگریزی \textit{joist}، فرانسیسی
\textit{solive}، اطالوی \textit{travicello}، روسی \textit{balka}، ہسپانوی تے
پرتگالی \textit{viga}۔ اک تکنیکی دھات جیہڑی ہالے اپݨائی نئیں گئی، پر ایتھے
فرانسیسی لفظ \textit{solive} دا ارتھ دیݨ لئی صلاح دِتی جاندی اے، جیہڑا نہ
\textit{trabo} (فرانسیسی \textit{poutre}) اے نہ نِکی کڑی۔ اوہ چوکور گھڑی لکڑ اے
جیہڑی اک فرش تے اک چھت سنبھالدی اے۔%
}

%%K t05-01-26 p
طویلے دی چھت \VUgras{کھپریلاں دی} \textsuperscript{(55)} اے، گھر دی
\VUgras{سلیٹ دی} \textsuperscript{(32)}، تے برامدے دی \VUgras{جست دی}
\textsuperscript{(24)}۔

%%K t05-01-27 p
\textsc{جان}۔ --- کیہ جست دیاں چھتاں وی چھت بݨاؤݨ والا ای بݨاندا اے؟

%%K t05-01-28 p
\textsc{مسٹر لبلاں}۔ --- نئیں، اوہ \VUgras{جست ساز} \textsuperscript{(121)} یا
\VUgras{نلساز} دا کم اے — اوہو ای جیہڑا گھر دی چھت وچ اک
\VUgras{روشندان کھڑکی} \textsuperscript{(38)} بہا رہیا اے۔ اوہو ای
\VUgras{چھت دے پرنالے} \textsuperscript{(43)} تے \VUgras{اُترن والیاں نالیاں}
\textsuperscript{(42)} وی لاندا اے۔ اوہ جست تے ٹین نوں ٹانکا لاندا اے۔ اوسے نے
\VUgras{بجلی روک} \textsuperscript{(40)} تے \VUgras{ہوا سُوچک} \textsuperscript{(39)}
\VUgras{مُنڈیر} \textsuperscript{(36)} اُتے گڈے سن۔

%%K t05-01-29 p
\textsc{جان}۔ --- تے کیہ گھر پورا ہو گیا؟

%%K t05-tit-5 sub
\VUsaut{3.0mm}
\VUpk{10.2pt}{40}{اوزار۔}
\VUsaut{3.0mm}

%%K t05-01-30 p
\textsc{مسٹر لبلاں}۔ --- پورا نئیں۔ \VUgras{پلستر کرن والا} \textsuperscript{(99)}
\VUgras{اِٹاں} \textsuperscript{(61)} نال اوہ کندھاں بݨاندا اے جیہڑیاں گھر نوں وکھو
وکھ کمریاں وچ ونڈدیاں نیں۔ اوہ \VUgras{پلستر} \textsuperscript{(70)}
\VUgras{چھتاں} \textsuperscript{(26)} تے کندھاں اُتے چاڑھدا اے۔ ایس ویلے اوہ
\VUgras{برامدے} \textsuperscript{(33)} دی چھت اُتے کم کر رہیا اے۔ راج وانگوں اوہدے
کول وی اک \VUgras{کرنی} \textsuperscript{(100)} تے اک \VUgras{نانْد}
\textsuperscript{(101)} اے۔

%%K t05-01-31 p
فیر ترکھاݨ دروازے، کھڑکیاں، \VUgras{لکڑ دے فرش} \textsuperscript{(16)} تے
\VUgras{جھلمِلاں} \textsuperscript{(19)} بݨاندا اے۔

%%K t05-01-32 p
ایس ویلے اوہ وِہڑے وچ اپݨی \VUgras{کاریگری دی میز} \textsuperscript{(90)} اُتے کم
کر رہیا اے۔ اوہدے کول اک \VUgras{آری} \textsuperscript{(94)} اے، \VUgras{لکڑ}
\textsuperscript{(69)} چیرن لئی، اک \VUgras{رندا} \textsuperscript{(91)}، اک
\VUgras{شکنجہ} \textsuperscript{(92)}، اک \VUgras{چھینی}
\textsuperscript{(94 \textit{bis})}، اک \VUgras{پرکار} \textsuperscript{(95)} تے
لکڑ دا اک \VUgras{مُنگلا} \textsuperscript{(95 \textit{bis})}۔

%%K t05-01-33 p
\textsc{جان}۔ --- اوہو، تدوں تے ترکھاݨ ہوݨا راج ہوݨ نالوں وی ودھ مزے دا اے۔ چاچا
جی، تُسیں مینوں اک کاریگری دی میز، اک آری تے اک رندا دلوا دیو گے؟ میں میدور لئی
اک گھر بݨاواں گا۔

%%K t05-01-34 p
\textsc{چاچا}۔ --- ہاں پُتر۔

%%K t05-01-35 p
\textsc{جان}۔ --- کِنّی خوشی ہوئی! تے اوہ بندہ کوݨ اے جیہڑا اگلے دروازے کول کھڑا
اے؟

%%K t05-tit-6 sub
\VUsaut{3.0mm}
\VUpk{10.2pt}{40}{رنگائی۔}
\VUsaut{3.0mm}

%%K t05-01-36 p
\textsc{مسٹر لبلاں}۔ --- اوہ \VUgras{رنگ ساز} \textsuperscript{(102)} اے۔ اوہ اپݨی
پوڑی اُتے کھڑا اے، \VUgras{رنگ} \textsuperscript{(103)} دی اک ہانڈی تے اپݨا
\VUgras{برش} \textsuperscript{(104)} چُکی۔ اوہ اگلے دروازے دی لکڑ اُتے رنگ چاڑھ
رہیا اے تاں جو اوہ ودھ ٹِکے۔

%%K t05-01-37 p
کمریاں وچ کندھ دا کاغذ وی اوہو ای چِپکاندا اے۔

%%K t05-01-38 p
\VUgras{سجاوٹ کرن والا} \textsuperscript{(107)}، جیہڑا \VUgras{پوڑی}
\textsuperscript{(106)} اُتے چڑھیا \VUgras{بالاخانے} \textsuperscript{(28)} اُتے اے،
اک \VUgras{پردہ} \textsuperscript{(29)} لا رہیا اے۔

%%K t05-01-39 p
وڈی کھڑکی کول کھڑا مزدور \VUgras{شیشہ ساز} \textsuperscript{(96)} اے۔ اوہ
\VUgras{شیشے} \textsuperscript{(18)} \VUgras{مصالحے} \textsuperscript{(73)} نال
جماندا اے۔ اوہ دوجا مزدور، جیہڑا تہہ خانے دے دروازے کول گوڈیاں بھار اے،
\VUgras{تالہ ساز} \textsuperscript{(97)} اے۔ اوہ اک \VUgras{تالا} \textsuperscript{(6)}
بہا رہیا اے۔ اوہدی \VUgras{ریتی} \textsuperscript{(98)} اوہدے کول پئی اے۔

%%K t05-01-40 p
\textsc{جان}۔ --- جیہڑا بندہ اپݨے \VUgras{ہتھوڑے} \textsuperscript{(84)} نال لوہے
دے ٹوٹے اُتے چوٹ کر رہیا اے، کیہ اوہ وی تالہ ساز اے؟

%%K t05-01-41 p
\textsc{مسٹر لبلاں}۔ --- ٹھیک آکھیے تے اوہ \VUgras{لوہار} \textsuperscript{(125)}
اے۔ اوہ \VUgras{لوہے دا سامان} \textsuperscript{(67)} گھڑ رہیا اے
\VUgras{بجلی والے} \textsuperscript{(124)} لئی، جیہڑا \VUgras{گڈی خانے}
\textsuperscript{(51)} دی چھت اُتے اے۔ اوہدے کول اک \VUgras{بھٹی} \textsuperscript{(126)}،
اک \VUgras{اہرݨ} \textsuperscript{(127)} تے اک \VUgras{سنڑاسی} \textsuperscript{(128)}
اے۔

%%K t05-01-42 p
بجلی والا ٹیلیفون دی \VUgras{تانبے دی تار} \textsuperscript{(44)} جوڑ رہیا اے۔

%%K t05-tit-7 sub
\VUpk{10.2pt}{40}{باغبانی۔}
\VUsaut{3.0mm}

%%K t05-01-43 p
\textsc{چاچا}۔ --- \VUgras{وِہڑے} \textsuperscript{(45)} دے سرے اُتے،
\VUgras{جنگلے} \textsuperscript{(47)} تے \VUgras{گڈی دے پھاٹک} \textsuperscript{(48)}
کول، تینوں \VUgras{مالی} \textsuperscript{(129)} وکھائی دیندا اے۔ اوہ نِکا
\VUgras{باغ} \textsuperscript{(49)} تیار کر رہیا اے۔ اوہنے اپݨی \VUgras{کہی}
\textsuperscript{(131)} نال مِٹی پُٹ سُٹی اے، اوہ بی بو رہیا اے تے \VUgras{پنجے}
\textsuperscript{(130)} نال پدھرا کر رہیا اے۔ فیر اپݨے \VUgras{جھارے}
\textsuperscript{(132)} نال اوہ \VUgras{کیاریاں} \textsuperscript{(150)} نوں سِنجے گا
تے بوٹے اُگݨ گے۔

%%K t05-01-44 p
\VUgras{سائیس} \textsuperscript{(133)}، جیہڑا ہُݨے گھوڑے دی سنبھال کر کے اوہنوں
کھوا چُکیا اے، \VUgras{بگھی} \textsuperscript{(52)} صاف کر رہیا اے — اسفنج، اک
\VUgras{ہتھ پمپ} \textsuperscript{(134)} تے پاݨی دی بالٹی نال۔ فیر اوہ اوہنوں گڈی
خانے وچ رکھ دیوے گا تے بھارے \VUgras{اڑگے} \textsuperscript{(53)} نال دروازہ بند کر
دیوے گا۔

%%K t05-01-45 p
لے، ہُݨ تے تسلی ہوئی ہوݨی؛ سمجھاؤݨا بہت ہو چُکیا۔

%%K t05-01-46 p
\textsc{جان}۔ --- ہاں چاچا جی، پر میں گھر دے اندر وی ویکھݨا چاہواں گا۔

%%K t05-01-47 p
\textsc{چاچا}۔ --- اوہ کل لئی۔ مسٹر رُو تینوں نقشہ سمجھاؤݨ گے تے ہر کمرے دا ناں
دسݨ گے۔

%%K t05-tit-8 sub
\VUsaut{3.0mm}
\VUpk{10.2pt}{40}{گھر دی اندرلی بݨتر۔}
\VUsaut{2.4mm}

%%K t05-apar-2 apar
{\centering\textit{(نقشہ ویکھو۔)}\par}
\VUsaut{2.4mm}

%%K t05-01-48 p
\textsc{معمار}۔ --- آ جان، میں تینوں گھر دے وکھو وکھ حصے وکھاندا واں۔

%%K t05-01-49 p
چل \VUgras{ہیٹھلی منزل} \textsuperscript{(7)} وچ چلیے۔ ایہہ رہیا \VUgras{گھنٹی} دا
بٹن \textsuperscript{(11)}۔ اسیں تن \VUgras{پوڑیاں} \textsuperscript{(14)} چڑھنے آں،
جیہڑیاں \VUgras{ڈیوڑھی} \textsuperscript{(9)} دیاں نیں، فیر \VUgras{دہلیز}
\textsuperscript{(10)} لنگھنے آں، جیہڑی \VUgras{اگلے بُوہے} \textsuperscript{(8)} دی
اے، تے ایہہ رہے اسیں \VUgras{زینے} \textsuperscript{(13)} دے تھلے۔

%%K t05-01-50 p
\textsc{جان}۔ --- ایہہ لانگھا اے۔

%%K t05-01-51 p
\textsc{معمار}۔ --- نئیں، ایہہ \VUgras{ڈیوڑھی کمرہ} \textsuperscript{(12)} اے۔
\VUgras{لانگھا} \textsuperscript{(a)} سرے اُتے اے۔ سجے پاسے ایہہ رہے
\VUgras{بیٹھک} \textsuperscript{(b)} تے \VUgras{کھاݨے دا کمرہ} \textsuperscript{(c)}؛
کھاݨے دے کمرے دے سامݨے \VUgras{رسوئی} \textsuperscript{(d)} تے
\VUgras{بھنڈار کمرہ} \textsuperscript{(e)}؛ فیر اگلے پاسے \VUgras{پڑھݨ دا کمرہ}
\textsuperscript{(f)}۔ \VUgras{برامدہ} \textsuperscript{(23)}، اپݨے
\VUgras{شیشے دے دروازے} \textsuperscript{(25)} سمیت، بیٹھک کول اے۔

%%K t05-01-52 p
ہُݨ اسیں زینہ چڑھاں گے۔ \VUgras{جنگلا} \textsuperscript{(15)} پھڑی رکھ تے پوڑیاں
گِݨ۔ چودہ نیں۔ ایہہ رہی \VUgras{چوکی} \textsuperscript{(m)}: اسیں پہلی
\VUgras{منزل} \textsuperscript{(27)} اُتے آں۔

%%K t05-tit-9 sub
\VUsaut{3.0mm}
\VUpk{10.2pt}{40}{پہلی منزل۔}
\VUsaut{3.0mm}

%%K t05-01-53 p
زینے دے سامݨے \VUgras{پاخانہ} \textsuperscript{(20)} تے \VUgras{کپڑیاں دا کمرہ}
\textsuperscript{(j)} نیں۔ سجے پاسے اک سوہݨا \VUgras{سوݨ دا کمرہ}
\textsuperscript{(g)}، جیدیاں دو کھڑکیاں گھر دے \VUgras{اگلے رُخ} \textsuperscript{(1)}
ول کھلدیاں نیں۔ کھبے پاسے \VUgras{نہاؤݨ دا کمرہ} \textsuperscript{(1)} تے
\VUgras{پراہُݨیاں دا کمرہ} \textsuperscript{(i)} نیں، جیدے وچ اک وڈا \VUgras{آلہ}
\textsuperscript{(h)} اے۔ سوݨ دے کمرے کول \VUgras{نِکیاں دا کمرہ}
\textsuperscript{(k)} اے۔ اوہ اک چوڑے \VUgras{چھت دے بارجے} \textsuperscript{(21)} وچ
کھلدا اے۔ ایس بارجے دے تھلے اک ہور پاخانہ \textsuperscript{(20)} اے، تے کندھ نال
لگی ایہہ رہی \VUgras{گھنٹی} \textsuperscript{(22)}، جیہڑی کھاݨے لئی وجائی جاندی اے۔

%%K t05-01-54 p
\textsc{جان}۔ --- چلو \VUgras{دوجی منزل} اُتے چلیے۔

%%K t05-01-55 p
\textsc{معمار}۔ --- دوجی منزل اے ای نئیں۔ چھت دے تھلے نِرے \VUgras{چُبارے}
\textsuperscript{(33)} تے بھنڈار \textsuperscript{(34)} نیں۔ ساڈا چکر پورا ہویا، ہُݨ
تھلے اُتر سکنے آں۔

%%K t05-01-56 p
\textsc{جان}۔ --- شکریہ جناب، بہت دلچسپ رہیا۔ میں اپݨے پیو نوں سب کجھ دساں گا
جیہڑا میں ویکھیا۔
\VUsaut{4.0mm}
\VUfilet{22mm}
